\chapter{Gleichungen und Ereignisse}
\label{kap:format-gleichungen}

Neben Schaltungen rechnet esim Gleichungssysteme, die direkt in der Datei
stehen – allein oder mit einer Schaltung gekoppelt. Damit lassen sich
Mechanik, Thermik, Reaktionskinetik, Regler oder eigene Bauteilmodelle
beschreiben.

\section{Variablen und Gleichungen}

\begin{json}
"variables": {
  "y1": 1,                               // Anfangswert
  "y2": { "init": 0, "abstol": 1e-12 },  // Anfangswert mit eigener Toleranz
  "y3": { "guess": 0.5 }                 // nur Startwert für den Löser
},
"equations": [
  "der(y1) = -0.04*y1 + 1e4*y2*y3",
  "der(y2) =  0.04*y1 - 1e4*y2*y3 - 3e7*y2^2",
  "0 = y1 + y2 + y3 - 1"
]
\end{json}

Es braucht genauso viele Gleichungen wie Variablen. Eine Gleichung ohne
\code{der()} ist eine algebraische Bedingung. \code{der()} darf einen
beliebigen Ausdruck der Unbekannten einschließen, muss aber linear mit
konstantem Faktor in der Gleichung stehen:

\begin{shell}
m*der(v) = F - k*x              ok
der(C0*u + C1*u^2) = -u/R       ok (nichtlineare "Ladung")
der(x) + der(y) = -x            ok
x*der(y) = 1                    nicht erlaubt -> Hilfsvariable einführen:
                                der(y) = dy  und  x*dy = 1
\end{shell}

Gleichungen höherer Ordnung werden als System erster Ordnung geschrieben
(\code{der(x) = v}, \code{der(v) = ...}). Variablen dürfen in Ausgängen,
Messgrößen, Verhaltensquellen und Schaltern verwendet werden.

\section{Anfangswerte}
\label{sec:anfangswerte}

Eine Variable, die in einem \code{der()} vorkommt, ist ein \emph{Zustand}.

\begin{itemize}
  \item Zahl oder \code{"init"}: Der Zustand startet bei diesem Wert.
  \item \code{"guess"} oder \code{null}: Der Zustand startet im
    Gleichgewicht (seine Zeitableitung ist bei $t=0$ null); der Wert ist nur
    Startwert für das Newton-Verfahren.
  \item Für algebraische Variablen ist jeder Wert nur ein Startwert: Vor
    der Integration werden sie so bestimmt, dass alle algebraischen
    Gleichungen erfüllt sind (konsistente Anfangswerte).
  \item \code{"abstol"} setzt die absolute Toleranz dieser Variable
    (Standard: \code{options.vartol}).
\end{itemize}

Die Analysen \code{op}, \code{dc} und \code{ac} rechnen immer mit dem
Gleichgewicht aller Gleichungen; die vorgegebenen Anfangswerte gelten nur
für \code{tran} und \code{pss}.

\section{Vektoren}
\label{sec:vektoren}

Viele gleichartige Unbekannte, etwa die Stützstellen einer örtlich
diskretisierten partiellen Differentialgleichung, werden als Vektor
vereinbart; eine Schleife schreibt ihre Gleichungen:

\begin{json}
"parameters": { "n": 41, "a": 1.1e-4, "dx": "0.2/(n-1)" },
"variables":  { "T[n]": 20 },
"equations": [
  "T[1] = 100",
  { "for": "i", "from": 2, "to": "n-1",
    "equation": "der(T[i]) = a * (T[i-1] - 2*T[i] + T[i+1]) / dx^2" },
  "T[n] = T[n-1]"
]
\end{json}

\begin{itemize}
  \item \code{"T[n]"} legt die Unbekannten \code{T[1]} bis \code{T[n]} an.
    Der Startwert ist eine Zahl für alle Elemente oder eine Liste mit einem
    Wert je Element (\code{"init": [...]}, \code{"guess": [...]}).
  \item Eine Schleife \code{{"for": "i", "from": .., "to": .., "equation": ".."}}
    steht für eine Gleichung je Indexwert. Mit \code{"equations": [...]}
    enthält sie mehrere Einträge, die wieder Schleifen sein dürfen.
  \item Der Index ist ein konstanter ganzzahliger Ausdruck aus Zahlen,
    Parametern und Schleifenindizes. Ein Index, der von einer Unbekannten
    abhängt, ist nicht möglich.
  \item \code{sum(i, von, bis, ausdruck)} summiert über einen Index.
  \item In den Ausgängen steht \code{T[*]} für alle Elemente.
\end{itemize}

Die Schleifen werden beim Einlesen ausgerollt; für den Löser sind es
einzelne Gleichungen. Matrix- und Vektoroperationen gibt es nicht; sie
lassen sich mit \code{sum} schreiben. Beispiel:
\datei{examples/heat_rod.json}.

\section{Indexreduktion}

Zwangsbedingungen der Mechanik und vorgegebene Bahnen führen auf Systeme,
die ein Integrationsverfahren nicht direkt lösen kann. Das Pendel in
kartesischen Koordinaten ist das Standardbeispiel:

\begin{json}
"equations": [
  "der(x) = u", "der(y) = v",
  "der(u) = -2*x*lam", "der(v) = -g - 2*y*lam",
  "0 = x^2 + y^2 - L^2"          // Index 3
]
\end{json}

Beim Einlesen bestimmt eine Strukturanalyse (Algorithmus von Pantelides),
wie oft jede Gleichung abgeleitet werden müsste. Ist der Index größer als
2, wird das System selbsttätig umgeformt:

\begin{itemize}
  \item Jede algebraische Bedingung $g = 0$, die $k$-mal abzuleiten wäre,
    erhält ihre Zeitableitungen bis zur Ordnung $k-1$ als zusätzliche
    Gleichungen. Sie werden symbolisch gebildet, wobei für \code{der(x)}
    die rechten Seiten der Gleichungen \code{c*der(x) = ...} eingesetzt
    werden.
  \item Je zusätzlicher Gleichung kommt ein Multiplikator
    \code{index_mu1}, \code{index_mu2}, \dots{} als Unbekannte hinzu. Er
    wirkt als $\dot x = \dots - \partial g/\partial x \cdot \mu$
    (Stabilisierung nach Gear, Gupta und Leimkuhler) und ist für die exakte
    Lösung null.
  \item Das Ergebnis hat Index 2. Bedingung und Ableitungen gelten mit der
    Genauigkeit der Newton-Iteration; es gibt keine Drift.
\end{itemize}

Das Programm meldet, was es hinzugefügt hat:

\begin{shell}
note: index 3 reduced to 2 by adding time derivatives of constraints:
  "0 = x^2 + y^2 - L^2":
    added  0 = 2*x*u + 2*y*v   (multiplier index_mu1 in the equations of x, y)
\end{shell}

\code{"options": {"index_reduction": false}} schaltet die Umformung ab.
Nicht reduzierbar sind Bedingungen, deren Ableitung sich nicht als Ausdruck
bilden lässt: wenn sie Knotenspannungen oder Ströme der Schaltung, Tabellen
oder Totzeiten enthalten, oder wenn für eine ihrer Unbekannten keine
Gleichung der Form \code{c*der(x) = ...} existiert. Dann bleibt das System,
wie es ist, und eine Warnung nennt die Ursache. Beispiele:
\datei{examples/pendulum.json} (von Hand auf Index 2 gebracht) und
\datei{examples/pendulum_index3.json} (automatisch).

\section{Ereignisse}
\label{sec:ereignisse}

\subsection*{Automatisch}

Jeder Vergleich (\code{<}, \code{<=}, \code{>}, \code{>=}) und jedes
\code{abs}, \code{sign}, \code{min}, \code{max} in Gleichungen und
Verhaltensquellen wird überwacht. Kippt eines davon während eines
Zeitschritts, sucht der Löser den Zeitpunkt, endet dort, schaltet um und
startet die Integration neu. Bis dahin rechnet er mit dem alten Zweig
weiter, sodass z.\,B. \code{der(z) = if(t < 1, 1, -2)} exakt bei $t = 1$
knickt. Dafür ist nichts weiter anzugeben. Schalter und ideale Dioden
schalten ebenso.

\subsection*{Mit Aktion}

Ein Eintrag unter \code{events} weist Zuständen neue Werte zu. Er hat
genau einen Auslöser:

\begin{center}
\begin{tabularx}{\textwidth}{@{}lL@{}}
  \toprule
  Auslöser & feuert \\
  \midrule
  \code{"when": "h < 0"} & in dem Moment, in dem die Bedingung wahr \emph{wird} \\
  \code{"every": "1m"} & periodisch; erstmals nach einer Periode oder zur Zeit \code{"start"} (z.\,B. \code{"start": 0}) \\
  \code{"at": 2.5} oder \code{"at": [1, 2.5]} & zu den angegebenen Zeiten \\
  \bottomrule
\end{tabularx}
\end{center}

\begin{json}
"variables": { "h": 1, "v": 0 },
"equations": [ "der(h) = v", "der(v) = -g" ],
"events": [
  { "when": "h < 0", "set": { "v": "-0.8*v" } }      // springender Ball
]
\end{json}

\begin{itemize}
  \item Die rechten Seiten von \code{set} werden mit den Werten vor dem
    Ereignis berechnet, erst dann wird zugewiesen. Löst eine Zuweisung
    weitere Bedingungen aus, feuern diese im selben Zeitpunkt.
  \item Zuweisen lassen sich Zustände, diskrete Variablen sowie Größen der
    Schaltung: \code{"v(c)"} für die Spannung eines Knotens mit
    Kondensator, \code{"i(L1)"} für einen Spulenstrom. Algebraische Größen
    ergeben sich nach dem Ereignis neu aus den Gleichungen; eine Zuweisung
    an sie bleibt wirkungslos.
  \item Zeitereignisse werden exakt getroffen. Mit \code{every} entsteht
    ein abgetastetes System: Ein Regler liest zu festen Zeiten die
    Schaltung aus und ändert seine diskreten Zustände, dazwischen läuft die
    Strecke kontinuierlich (\datei{examples/digital_pi_control.json}).
  \item Eine Bedingung, die schon beim Start wahr ist, löst nicht aus.
\end{itemize}

\subsection*{Diskrete Variablen}

Diskrete Variablen ändern sich nur durch Zuweisungen und sind dazwischen
konstant: Schaltzustände, Betriebsarten, Zähler. Sie brauchen keine
Gleichung:

\begin{json}
"variables": { "T": 15, "heat": { "discrete": 1 } },
"equations": [ "Cth*der(T) = P*heat - (T - Tout)/Rth" ],
"events": [
  { "when": "T > 22", "set": { "heat": 0 } },         // Zweipunktregler
  { "when": "T < 20", "set": { "heat": 1 } }          // mit Hysterese
]
\end{json}

In \code{op}, \code{dc} und \code{ac} behalten diskrete Variablen ihren
Startwert.

\subsection*{Grenzen}

Ein Modell, das ohne Hysterese hin- und herschaltet
(\code{der(x) = if(x < 1, 1, -1)} pendelt um $x = 1$) oder dessen
Ereignisse sich häufen (der Ball springt unendlich oft in endlicher Zeit),
bricht nach \code{max_events} Ereignissen mit einer Fehlermeldung ab.
Solche Modelle brauchen eine Hysterese oder eine geglättete Kennlinie. Der
Löser sucht Vorzeichenwechsel am Ende und an drei Stellen im Inneren jedes
Zeitschritts; ein Hin- und Rückwechsel, der kürzer als ein Viertelschritt
ist, kann unbemerkt bleiben. \code{tstep} begrenzt die Schrittweite.

\section{Totzeiten}
\label{sec:totzeit}

\code{delay(x, td)} liefert den Wert, den der Ausdruck \code{x} vor der
konstanten Zeit \code{td} hatte:

\begin{json}
"equations": [ "T * der(y) = K * delay(u, Td) - y" ]
\end{json}

\begin{itemize}
  \item Vor dem Start gilt der Wert bei $t=0$ (das System war in Ruhe).
  \item Im Arbeitspunkt und in der AC-Analyse ist \code{delay(x, td)}
    gleich \code{x}.
  \item Kein Zeitschritt ist länger als die kleinste Totzeit. Zwischen den
    gespeicherten Zeitpunkten wird kubisch interpoliert, an Knicken linear;
    Knicke und Sprünge setzen eine Totzeit später einen Haltepunkt.
  \item \code{td} muss eine positive Konstante sein (Zahl oder Parameter).
\end{itemize}

\datei{examples/dead_time_control.json} zeigt einen Regelkreis mit
Totzeit.

\section{Kopplung mit der Schaltung}

Gleichungen lesen die Schaltung über \code{v(...)} und \code{i(...)}. In die
andere Richtung wirkt die Verhaltensquelle \code{B}, deren Strom, Spannung
oder Ladung ein Ausdruck ist, der auch Variablen enthalten darf:

\begin{json}
"variables": { "T": 298.15, "Rntc": { "guess": 10 } },
"equations": [
  "Cth*der(T) = v(in,mid)^2/Rntc - (T - Tamb)/Rth",     // Wärmebilanz
  "Rntc = R25*exp(B*(1/T - 1/T0))"                       // Kennlinie
],
"components": [
  { "type": "B", "name": "NTC", "nodes": ["in", "mid"], "i": "v(in,mid)/Rntc" }
]
\end{json}

Das vollständige Beispiel ist \datei{examples/ntc_inrush.json}: ein
Heißleiter, der sich durch seinen eigenen Strom erwärmt und den
Einschaltstrom begrenzt. In eine Teilschaltung verpackt, wird daraus ein
wiederverwendbares Bauteil mit thermischem Gedächtnis.
